\subsection{Noise}\label{sec:theory_general_noise}
Any unwanted signal that is present in a circuit is called noise. It is
usually measured in terms of its power spectral density (PSD) and then
converted to root mean square (RMS) voltage or current. Some common noises
are:
\begin{itemize}
  \item Thermal Noise: This is the noise generated by the thermal fluctuations
    of charge carriers (usually electrons) in a conductor or semiconductor.
    It is also called Johnson-Nyquist noise, is given by:
    \begin{align*}
      S_v(f) &= 4 k T R \hspace{1cm} \text{[in V\(^2\)/Hz; voltage PSD, source in series with \(R\)]} \\
      S_i(f) &= \frac{4 k T}{R} = 4kTG \hspace{1cm} \text{[in A\(^2\)/Hz; current PSD, source in parallel]} \\
      V_n &= \sqrt{\int_{f_L}^{f_H} 4 k T R \, df} = \sqrt{4 k T R (f_H - f_L)} \hspace{1cm} \text{[in V\(_{\text{rms}}\)]}
    \end{align*}
    where \(k\) is Boltzmann's constant, \(T\) is the absolute temperature in
    Kelvin, \(R\) is the resistance in ohms, and \(f_H - f_L\) is the
    bandwidth of the signal. A resistor in parallel with a capacitor
    integrates to the bandwidth independent value
    \(\overline{V_n^2} = \int_0^\infty \frac{4kTR}{1 + (2\pi f R C)^2} df = \frac{kT}{C}\).
  \item Shot Noise: This is the noise generated by the discrete and probabilistic
    nature of charge carriers (usually electrons) crossing a potential barrier
    (a PN junction). It is given by:
    \begin{align*}
      S_i(f) &= 2 q I \hspace{1cm} \text{[in A\(^2\)/Hz; current PSD]} \\
      I_n &= \sqrt{2 q I (f_H - f_L)} \hspace{1cm} \text{[in A\(_{\text{rms}}\)]}
    \end{align*}
    where \(q\) is the charge of an electron and \(I\) is the average current
    in amperes.
  \item Flicker Noise: This is the noise that has a power spectral density that
    is inversely proportional to frequency. It is also called 1/f noise and
    is given by:
    \begin{align*}
      S_v(f) &= \frac{K}{f^a} \hspace{1cm} \text{[in V\(^2\)/Hz; voltage PSD]}
    \end{align*}
    where \(K\) is a constant that depends on the device and its operating
    conditions and \(a\) is an exponent that typically ranges from 0.5 to 2.
    Since the PSD is not flat, it must be integrated over the band rather
    than multiplied by it:
    \begin{align*}
      V_n^2 &= \int_{f_L}^{f_H} \frac{K}{f^a} \, df =
      \begin{cases}
        K \ln\left(\frac{f_H}{f_L}\right) & a = 1 \\[2mm]
        \frac{K}{1 - a}\left(f_H^{1-a} - f_L^{1-a}\right) & a \neq 1
      \end{cases}
    \end{align*}
    For \(a = 1\) every decade of bandwidth adds the same noise power \(K \ln 10\).
  \item Uncorrelated sources add in power:
    \(V_{n,total} = \sqrt{V_{n1}^2 + V_{n2}^2 + \cdots}\).
\end{itemize}
